% ID: SEL-D0004-OA-E03
% ESTADO: validada. S005, 2026-11-07; tres pasadas del mismo autor y revisión matemática, didáctica y visual.
\textbf{Datos y partición.}\par\nopagebreak[4]
Sea $S$ ``superar Matemáticas''. Los institutos forman una partición:
\[
P(A)=250/500=0.5,\quad P(B)=150/500=0.3,\quad P(C)=100/500=0.2.
\]
Las tasas de aprobado son $P(S\mid A)=0.70$, $P(S\mid B)=0.68$ y $P(S\mid C)=0.73$.

\textbf{a) Probabilidad total de no superar la materia.}\par\nopagebreak[4]
\[
\begin{aligned}
P(S^c)&=0.5(0.30)+0.3(0.32)+0.2(0.27)\\
&=0.15+0.096+0.054=\boxed{0.30}.
\end{aligned}
\]
No se hace una media simple de porcentajes: los tamaños de los institutos son distintos.

\textbf{b) Instituto A sabiendo que ha aprobado: Bayes.}\par\nopagebreak[4]
Como $P(S)=1-0.30=0.70$ y $P(A\cap S)=0.5(0.70)=0.35$,
\[
P(A\mid S)=\frac{P(A\cap S)}{P(S)}=\frac{0.35}{0.70}=\boxed{0.5}.
\]
La tasa dada $P(S\mid A)$ no es la probabilidad pedida $P(A\mid S)$.

\textbf{c) Instituto C y no superar la materia.}\par\nopagebreak[4]
\[
P(C\cap S^c)=P(C)P(S^c\mid C)=0.2(0.27)=\boxed{0.054}.
\]
Comprobación por recuentos: aprueban $175+102+73=350$ de los 500; suspenden 150. De C suspenden 27, por lo que $27/500=0.054$; entre aprobados, $175/350=0.5$ son de A.
